% Versione aggiornata e tradotta in LaTeX di
% documentazione_soldani/2_1_benchmark_thor_description.txt (12/07/2026),
% ormai obsoleto su piu' punti. Aggiornato al 20/07/2026 dopo l'evoluzione
% del notebook, l'aggiunta del modulo src_soldani, della pipeline
% multimodale e il debug dell'infrastruttura Thor.
% Richiede \usepackage{booktabs} e \usepackage{amsmath} nel preambolo.

\section{Descrizione del benchmark FairMind vs LLM su THOR}

\subsection*{Nota di aggiornamento}
Questa sezione sostituisce e aggiorna \texttt{2\_1\_benchmark\_thor\_description.txt} (12 luglio 2026), che descriveva un notebook nel frattempo rinominato e un'infrastruttura Thor mai effettivamente verificata end-to-end. Le modifiche principali rispetto alla versione originale sono segnalate esplicitamente in ciascuna sottosezione.

\subsection{Introduzione e obiettivo}

Il benchmark confronta due approcci per il calcolo degli effetti di causal fairness sul dataset Adult Income (UCI): \emph{FairMind}, la libreria del progetto che costruisce uno Standard Fairness Model (SFM, Ple\v{c}ko e Bareinboim, 2024) sotto forma di Bayesian Network e calcola gli effetti con inferenza causale formale; e un \emph{LLM locale} (Qwen, servito via \texttt{llama.cpp} sul cluster HPC THOR), a cui vengono fornite tabelle di probabilit\`a pre-aggregate e le formule di identificazione, chiedendogli di applicarle.

\textbf{Aggiornamento rispetto alla versione originale:} il lavoro copre ora due notebook distinti, non pi\`u uno solo:
\begin{itemize}
  \item \texttt{2\_3\_benchmark\_thor.ipynb} --- benchmark singolo (ex \texttt{2\_1\_benchmark\_thor.ipynb}, rinominato nel commit \texttt{9391039});
  \item \texttt{3\_causal\_analysis\_pipeline.ipynb} --- pipeline completa aggiunta successivamente, che estende il confronto con causal discovery automatica e analisi intersezionale (si veda \S\ref{sec:pipeline}).
\end{itemize}

I cinque effetti calcolati restano invariati:
\[
\text{TV},\ \text{TE} = \text{TV} - \text{SE},\ \text{SE},\ \text{DE},\ \text{IE}
\]
rispettivamente Total Variation, Total Effect, Spurious Effect, Natural Direct Effect e Natural Indirect Effect.

\subsection{Contesto: il cluster THOR}

THOR \`e il cluster HPC della SUPSI (Slurm 25, Apptainer 1.4.2), accessibile via \texttt{ssh} a \texttt{thor.supsi.ch}.

\textbf{Aggiornamento maggiore:} la descrizione originale della configurazione Thor era in gran parte teorica/mai testata. Il debug del 20 luglio 2026 (si veda \S\ref{sec:debug-thor} e il diario di bordo) ha rivelato diverse discrepanze tra quanto documentato e quanto realmente funzionante:

\begin{itemize}
  \item Il server LLM persistente (job Slurm \texttt{llama\_server\_gpu}, non \texttt{llama\_server} come nominato in origine) gira sulla partition \texttt{gpu} (non \texttt{gpuISIN}: l'utente non ha i permessi per quest'ultima) e serve il modello \texttt{Qwen2.5-14B-Instruct-Q4\_K\_M.gguf}. Il codice del client (\texttt{src/llm.py}, campo \texttt{model="qwen2.5-7b-instruct"}) dichiara invece un 7B: il nome del modello nella richiesta API non seleziona il file effettivamente servito (il server ne serve uno solo), quindi le risposte finora ottenute --- inclusi tutti i confronti storici nel report --- provengono verosimilmente dal modello 14B, non dal 7B come sempre dichiarato in documentazione precedente. Punto da chiarire con il relatore prima di riportare risultati definitivi sul modello utilizzato.
  \item I container Apptainer usati realmente sono \texttt{containers/llama\_server\_gpu.sif} e \texttt{containers/jupyter\_env.sif} (percorso \texttt{scratch/causal\_fairness/containers/}, non la root di \texttt{scratch/} come indicato in precedenza).
  \item I job che eseguono i notebook (\texttt{run\_benchmark.sbatch}, \texttt{run\_multimodal.sbatch}) non avviano pi\`u un proprio \texttt{llama-server}: si agganciano a quello gi\`a attivo cercandolo per nome job via \texttt{squeue}. In precedenza ciascuno ne avviava uno proprio, richiedendo una seconda GPU e restando spesso bloccato indefinitamente in \texttt{PENDING}.
  \item La partition \texttt{compute} presenta un problema di risoluzione utente (NSS/LDAP) che fa fallire silenziosamente sia \texttt{squeue -u \$USER} sia \texttt{apptainer exec} all'interno di un job batch; i notebook vengono quindi eseguiti sulla partition \texttt{gpu} senza richiedere risorse GPU (\texttt{--gres} omesso).
  \item Le dipendenze Python dentro \texttt{jupyter\_env.sif} sono ora pinnate alle stesse versioni del \texttt{.venv} locale (\texttt{pandas==2.3.3}, \texttt{pgmpy==1.0.0}, ecc.): versioni pi\`u recenti scaricate senza pin avevano rotto sia la firma di \texttt{DiscreteBayesianNetwork.fit()} sia l'inferenza automatica del tipo di dato per colonne stringa.
\end{itemize}

\subsection{Evoluzione dei notebook di benchmark}

\begin{center}
\begin{tabular}{lll}
\toprule
Notebook & Backend LLM & Note \\
\midrule
\texttt{2\_1\_benchmark\_gpt.ipynb} & OpenAI (\texttt{o3-mini}) & dati grezzi CSV nel prompt, richiede internet \\
\texttt{2\_3\_benchmark\_thor.ipynb} & \texttt{llama.cpp} locale su Thor & tabelle pre-aggregate, ex \texttt{2\_1\_benchmark\_thor} \\
\texttt{3\_causal\_analysis\_pipeline.ipynb} & \texttt{llama.cpp} locale su Thor & discovery + intersezionale + report HTML \\
\bottomrule
\end{tabular}
\end{center}

La scelta delle tabelle pre-aggregate resta motivata dalla finestra di contesto limitata di Qwen2.5-7B/14B: inviare le \textasciitilde 48\,842 righe del dataset Adult come CSV grezzo (approccio \texttt{2\_1\_benchmark\_gpt}) produce un prompt impraticabile e, come mostrato dal primo benchmark, un LLM locale che restituisce output nulli dopo oltre 10 minuti.

\subsection{Setup e connessione al LLM}

\textbf{Aggiornato:} la cella di configurazione del client LLM non contiene pi\`u un'istanza \texttt{OpenAI} locale hardcoded n\'e la variabile \texttt{MODEL\_NAME} (entrambe codice morto, rimosso il 20/07/2026). La connessione avviene ora tramite \texttt{src.llm.LLM\_CONFIGS}, la cui voce di default legge le variabili d'ambiente \texttt{LLAMA\_HOST}/\texttt{LLAMA\_PORT} (fallback \texttt{localhost:8080}):

\begin{verbatim}
LLAMA_HOST = os.environ.get("LLAMA_HOST", "localhost")
LLAMA_PORT = os.environ.get("LLAMA_PORT", "8080")
LLM_CONFIGS[0]["base_url"] = f"http://{LLAMA_HOST}:{LLAMA_PORT}/v1"
\end{verbatim}

Questo permette allo stesso notebook di funzionare sia in locale sia su Thor, agganciandosi al nodo corretto senza modifiche manuali al codice. \texttt{src/llm.py} include inoltre, rispetto alla versione originale, il supporto multi-modello (\texttt{register\_llm}, \texttt{clear\_llm\_configs}, \texttt{benchmark\_models}) usato dalla pipeline (\S\ref{sec:pipeline}) per confrontare pi\`u LLM sullo stesso prompt.

\subsection{Configurazione del benchmark}

Invariata rispetto alla versione originale: dataset Adult, target \texttt{T\_income} (\texttt{>50K}), attributo protetto \texttt{S2\_gender} ($x_0{=}$Female, $x_1{=}$Male), mediatore \texttt{hours-per-week}, confondente \texttt{education}.

\subsection{FairMind --- ground truth (con correzione del binning)}
\label{sec:binning-fix}

\textbf{Bug corretto il 20/07/2026, non presente nella descrizione originale.} La versione precedente affermava che ``il dataset Adult preprocessato ha gi\`a tutte le variabili categoriche, quindi non serve discretizzazione esplicita'': questa affermazione era \emph{errata} per il mediatore \texttt{hours-per-week}, che \texttt{run\_fairmind()} riceveva nella sua forma continua originale, mentre il prompt per il LLM lo discretizzava gi\`a in 5 fasce (\texttt{<=20}, \texttt{21-35}, \texttt{36-45}, \texttt{46-60}, \texttt{>60}). I due metodi confrontavano quindi effetti calcolati su rappresentazioni diverse della stessa variabile per DE e IE, gli unici due effetti che dipendono dal mediatore.

Corretto applicando la stessa discretizzazione anche a monte del calcolo FairMind. Il valore del ground truth per DE e IE \`e conseguentemente cambiato (si veda la tabella aggiornata in \S\ref{sec:risultati}); TV, TE e SE, che non dipendono dal mediatore, restano identici.

\subsection{Prompt engineering per il LLM}

Invariato: le cinque tabelle di probabilit\`a condizionale (P(Y|X), P(Z), P(Y|X,Z), P(W|X,Z), P(Y|X,W,Z)) vengono calcolate con \texttt{groupby} di pandas e inserite nel prompt insieme alle formule di identificazione, per un totale di circa 11\,000 caratteri (\textasciitilde 7000 token).

\subsection{Chiamata al LLM e parsing della risposta}

\textbf{Aggiornato:} la funzione non si chiama pi\`u \texttt{call\_gpt} ma \texttt{call\_llm} (rinominata nel commit \texttt{9391039} insieme a tutti i riferimenti \texttt{gpt}\,$\to$\,\texttt{llm} nel codice e nelle celle markdown). La logica di parsing (ricerca di blocchi JSON con regex, filtraggio per le 5 chiavi attese, presa dell'ultimo blocco valido) \`e invariata.

\subsection{Confronto FairMind vs LLM: risultati aggiornati}
\label{sec:risultati}

\begin{center}
\begin{tabular}{lrrrr}
\toprule
Effetto & FairMind (12/07, non corretto) & FairMind (20/07, corretto) & LLM & Errore relativo \\
\midrule
TV & 0.194470 & 0.194470 & 0.194500 & 0.02\% \\
TE & 0.183161 & 0.183161 & 0.231167 & 26.21\% \\
SE & -0.007296 & -0.007296 & -0.036667 & 402.56\% \\
DE & 0.137049 & \textbf{0.137359} & 0.083167 & 39.45\% \\
IE & -0.046112 & \textbf{-0.045803} & -0.047000 & 2.61\% \\
\bottomrule
\end{tabular}
\end{center}

I valori del LLM sono invariati (la sua discretizzazione era gi\`a corretta); solo il ground truth FairMind cambia per DE/IE, con un piccolo ma reale impatto sull'errore relativo riportato per questi due effetti. Le conclusioni qualitative restano le stesse: il LLM \`e molto accurato su TV e IE, mediocre su TE/DE, inaffidabile su SE (per l'amplificazione dell'errore relativo su un valore vero prossimo a zero).

\subsection{Salvataggio risultati}

Invariato nella struttura (\texttt{benchmark\_results/\{dataset\}\_\{timestamp\}.json}), con la sola differenza che la chiave del payload relativa al LLM si chiama ora \texttt{llm} anzich\'e \texttt{gpt}, coerentemente con il rename di \S 9.

\subsection{Modulo \texttt{src\_soldani/} (nuovo)}

Non presente nella descrizione originale: modulo aggiunto per estendere FairMind oltre il confronto singolo con il LLM.

\begin{itemize}
  \item \texttt{discovery.py} --- \texttt{discover\_graph}/\texttt{learn\_sfm}: apprendimento strutturale del grafo causale (Hill-Climb o PC algorithm) invece di specificarlo manualmente; \texttt{graph\_similarity}: confronto tra due grafi (SHD, precision, recall, F1, Jaccard).
  \item \texttt{intersectional.py} --- \texttt{combine\_sensitive\_attrs}/\texttt{compute\_intersectional\_effects}: calcolo dei 5 effetti per ogni coppia di gruppi intersezionali (es. genere $\times$ razza).
  \item \texttt{report.py} --- \texttt{generate\_html\_report}: report HTML con tutte le sezioni del confronto (configurazione, grafo scoperto vs manuale, ground truth, benchmark multi-modello, discrepanze, analisi intersezionale, timing).
\end{itemize}

\subsection{Pipeline multimodale: \texttt{3\_causal\_analysis\_pipeline.ipynb} (nuovo)}
\label{sec:pipeline}

Notebook aggiunto dopo la stesura della descrizione originale, assente da quella versione. Esegue in sequenza:
\begin{enumerate}
  \item \textbf{Causal discovery} (Hill-Climb, score BDeu) sul grafo appreso dai dati (con \texttt{hours-per-week} gi\`a discretizzato, cos\`i come nel prompt LLM);
  \item \textbf{FairMind su grafo manuale} (ground truth di riferimento, stesso fix del binning di \S\ref{sec:binning-fix});
  \item \textbf{FairMind su grafo scoperto} --- qui \`e stato corretto un secondo caso dello stesso bug: il fitting del modello bayesiano avveniva su dati grezzi mentre la struttura era stata appresa su dati gi\`a binnati, disallineamento pi\`u grave del precedente perch\'e non solo numerico ma strutturale;
  \item \textbf{Analisi intersezionale} (genere $\times$ razza, tutte le coppie);
  \item \textbf{Costruzione del prompt LLM};
  \item \textbf{Benchmark multi-modello} (\texttt{benchmark\_models}, supporta pi\`u LLM sullo stesso prompt);
  \item \textbf{Generazione del report HTML} (\texttt{generate\_html\_report}).
\end{enumerate}
Eseguito da \texttt{run\_multimodal.sbatch}, aggiornato lo stesso giorno del debug con le stesse correzioni infrastrutturali di \texttt{run\_benchmark.sbatch}.

\subsection{Debug dell'infrastruttura Thor (20 luglio 2026)}
\label{sec:debug-thor}

Sintesi dei problemi trovati e risolti in una singola sessione di debug, descritti in dettaglio nella voce di diario \texttt{giorni/20260720.tex}:
\begin{enumerate}
  \item Nome del job cercato (\texttt{llama\_server}) non corrispondente a quello reale (\texttt{llama\_server\_gpu}).
  \item Permessi mancanti per la partition \texttt{gpuISIN}, causa di job bloccati indefinitamente in \texttt{PENDING}.
  \item Risoluzione utente NSS/LDAP rotta sui nodi \texttt{compute} all'interno di job batch (\texttt{squeue}, \texttt{whoami}, \texttt{apptainer exec} falliscono tutti con lo stesso errore di fondo).
  \item Controllo dell'exit code mancante: i job riportavano falso successo anche quando \texttt{apptainer}/\texttt{jupyter nbconvert} fallivano.
  \item File \texttt{.def} non validi (\texttt{llama\_server.def} incollato dentro un heredoc bash) o con typo (\texttt{jupyter\_env.def}, \texttt{--no-cache-cdir}).
  \item Dipendenze Python mancanti (\texttt{plotly}, \texttt{daft-pgm}) o non pinnate (\texttt{pgmpy}, \texttt{pandas}), causa di due rotture successive dell'esecuzione.
\end{enumerate}
Al termine della sessione, \texttt{run\_benchmark.sbatch} ha completato con successo l'intera pipeline in circa 44 secondi, primo run end-to-end riuscito su Thor.

\subsection{Conclusioni aggiornate}

Le conclusioni originali restano valide nella sostanza: FairMind \`e il riferimento per accuratezza e velocit\`a (\textasciitilde 3500 volte pi\`u veloce), il LLM \`e affidabile solo per gli effetti pi\`u semplici (TV, IE). A questo si aggiungono, rispetto alla versione del 12 luglio, due considerazioni nuove:
\begin{itemize}
  \item la correttezza del confronto dipende in modo sottile dall'allineamento tra le rappresentazioni delle variabili usate da FairMind e dal LLM (si veda il bug del binning, \S\ref{sec:binning-fix}) --- un aspetto metodologico non ovvio da verificare a occhio;
  \item l'esecuzione end-to-end su un cluster HPC condiviso introduce una classe di problemi (permessi, risoluzione utente, versioning delle dipendenze) indipendente dalla correttezza del codice applicativo, e va documentata con la stessa cura riservata al codice stesso.
\end{itemize}
